%%% Local Variables:
%%% mode: latex
%%% TeX-master: t
%%% End:
% The following command is used with glossaries-extra
\setabbreviationstyle[acronym]{long-short}
% The form of the entries in this file is \newacronym{label}{acronym}{phrase}
%                                      or \newacronym[options]{label}{acronym}{phrase}
% see "User Manual for glossaries.sty" for the  details about the options, one example is shown below
% note the specification of the long form plural in the line below
%
% The following example also uses options

% note the use of a non-breaking dash in long text for the following acronym

% example of putting in a trademark on first expansion


% note the use of a non-breaking dash in the following acronym


% ---- Acronyms used in this thesis ----
\newacronym{ACR}{ACR}{absolute category rating}
\newacronym{GNLL}{GNLL}{Gaussian negative log-likelihood}
\newacronym{ITU-T}{ITU-T}{ITU Telecommunication Standardization Sector}
\newacronym{MOS}{MOS}{mean opinion score}
\newacronym{NaN}{NaN}{not a number}
\newacronym{PCC}{PCC}{Pearson correlation coefficient}
\newacronym{PESQ}{PESQ}{perceptual evaluation of speech quality}
\newacronym{POLQA}{POLQA}{perceptual objective listening quality assessment}
\newacronym{PSTN}{PSTN}{public switched telephone network}
\newacronym{ReLU}{ReLU}{rectified linear unit}
\newacronym{RMSE}{RMSE}{root-mean-square error}
\newacronym{SSL}{SSL}{self-supervised learning}
\newacronym{TTS}{TTS}{text-to-speech}
\newacronym{VoIP}{VoIP}{voice over Internet Protocol}
